% generated by roam.py handbook export -- не редактировать, править заметки
% 20260920221509-матрица_гессе.org (body 7019915fa82d)

\begin{definition}[{Матрица Гессе}]\label[definition]{note:592AD816-1D80-48E6-8399-8524FBC1E8F3}
  Пусть $E \subset \R^n$~--- открытое множество, $f: E \rightarrow \R$~--- дважды непрерывно дифференцируемая функция.
  \textbf{Матрица Гессе} (\emph{гессиан}) $f$ в точке $x \in E$~--- матрица $n \times n$ из вторых частных производных:
  \[
  \nabla^2 f(x) = \left( \frac{\partial^2 f}{\partial x_i \partial x_j}(x) \right)_{i,j=1}^n
  \]

  Смешанные производные непрерывны, поэтому не зависят от порядка дифференцирования, и $\nabla^2 f(x)$~--- симметричная матрица.

  Матрица Гессе~--- это производная \hyperref[note:1756907B-1FF2-4F62-8553-759B8509281C]{градиента} как отображения $\nabla f: E \rightarrow \R^n$: ее $i$-я строка состоит из частных производных $i$-й координаты градиента $\frac{\partial f}{\partial x_i}$.

  Запись $\nabla^2 f(x) \ge 0$ означает \emph{неотрицательную определенность}: $h^T \nabla^2 f(x) h \ge 0\ \forall h \in \R^n$. Это нестрогий аналог \hyperref[note:771B93FA-654D-4233-9487-B4FDD356DDFE]{положительной определенности}, где требуется $h^T A h > 0\ \forall h \ne 0$.
\end{definition}

\begin{proposition}[{Вторая производная вдоль прямой}]\label[proposition]{note:592AD816-1D80-48E6-8399-8524FBC1E8F3:1}
  Пусть $x \in E$, $h \in \R^n$ и $\phi(t) = f(x + th)$ (для тех $t$, при которых $x + th \in E$). Тогда
  \[
  \phi^{\prime\prime}(t) = \left\langle \nabla^2 f(x + th) h; h \right\rangle = h^T \nabla^2 f(x + th) h
  \]
\end{proposition}

\begin{proof}
  Из заметки про \hyperref[note:1756907B-1FF2-4F62-8553-759B8509281C]{градиент} (производная вдоль прямой):
  \[
  \phi^{\prime}(t) = \left\langle \nabla f(x + th); h \right\rangle = \sum\limits_{i=1}^n \frac{\partial f}{\partial x_i}(x + th) h_i
  \]
  Каждое слагаемое~--- снова функция вида $u(x + th)$, где $u = \frac{\partial f}{\partial x_i}$, поэтому к нему применима та же формула:
  \[
  \left( \frac{\partial f}{\partial x_i}(x + th) \right)^{\prime}_t = \left\langle \nabla \frac{\partial f}{\partial x_i}(x + th); h \right\rangle = \sum\limits_{j=1}^n \frac{\partial^2 f}{\partial x_i \partial x_j}(x + th) h_j
  \]
  Подставим:
  \[
  \phi^{\prime\prime}(t) = \sum\limits_{i=1}^n \sum\limits_{j=1}^n \frac{\partial^2 f}{\partial x_i \partial x_j}(x + th) h_i h_j = h^T \nabla^2 f(x + th) h
  \]
\end{proof}
